\documentclass{article}
\usepackage{tcee-repaso}
\tcetema[Fiscal, crecimiento y ahorro]{3.A.38}{La política fiscal: efectos sobre el crecimiento económico y el ahorro.}
\tcefecha{05/11/2026}
\begin{document}
\cabecerarepaso
\begin{columnas}

% ==========================================================================
\bloque{Esquema del cante}
\esquemacante{%
\apartadocante{Introducción: la cigarra y la hormiga}{2,5}
\apartadocante{I. Corto plazo: estabilizar la demanda}{9,5}
\subcante{1. Debate tradicional: IS-LM}{1}
\subcante{2. Modelo de ciclo real con gasto}{5,5}
\subcante{3. Modelo neokeynesiano}{3}
\apartadocante{II. Largo plazo: el ahorro}{9}
\subcante{1. Agente de horizonte infinito}{3}
\subcante{2. Pensiones públicas y ahorro}{6}
\apartadocante{III. Largo plazo: el crecimiento}{8}
\subcante{1. Crecimiento endógeno: \autor{Barro} (1990)}{5,5}
\subcante{2. Estructura impositiva y crecimiento}{2,5}
\apartadocante{Conclusión}{1}}
\emph{Reparto}: el título pide crecimiento y ahorro: corto plazo $\leq$\tiempo{9,5}; IS-LM en \tiempo{1}; NEK y multiplicadores, resumidos \vertema{3.A.40} \vertema{4.B.4}. 3 desarrollos propios: ciclo real con gasto y ER, pensiones en \autor{Diamond} y \autor{Barro} (1990). \emph{Hilo}: \emph{nivel} frente a \emph{tasa}. \emph{Funciones clave}: recursos con $G$, \autor{Euler} con $\tau_K$, ahorro con reparto, IS con gasto y $\tau^*=1-\alpha$.

% ==========================================================================
\bloque[2,5]{Introducción}
\begin{itemize}
\item \textbf{Enganche}: «Quien no trabaja en verano, pasa hambre en invierno»: la fábula de \autor{Samaniego} (1781), la cigarra canta, la hormiga guarda. Pero la hormiga también gasta esfuerzo en el hormiguero: no ahorrar no es ser cigarra si se gasta en lo que rinde mañana (infraestructuras, formación, I+D). Y el reparto es un granero común: el invierno de cada generación depende del verano de la siguiente (pensiones, 13,1\,\% del PIB en 2022 $\to$ máx.\ 17,3\,\% hacia 2050). ¿\emph{Puede la política fiscal hacer de hormiga, gastando y ahorrando para crecer más, o solo adelanta el verano a costa del invierno}?
\item \textbf{Relevancia}: \emph{teoría}: se cruzan ciclo y crecimiento; los efectos dependen del horizonte de los hogares, de la financiación (suma fija, distorsionantes, deuda) y de si el gasto es productivo. \emph{Política}: gasto AA.\,PP.\ 45,3\,\% del PIB en España y 49,8\,\% en la zona del euro (2025); reglas UE de 2024: senda de gasto neto \vertema{3.A.39}; envejecimiento \vertema{4.B.24}. \emph{España 2025}: déficit 2,4\,\%, deuda 100,5\,\%; ahorro nacional bruto 25,0\,\% (máx.\ desde 1995) y AA.\,PP.\ vuelven a ahorrar (0,3\,\%; no desde 2008): ahorro público = ingresos $-$ gasto corriente.
\item \textbf{Contextualización}: \concepto{política fiscal} = actuación del sector público sobre gasto, impuestos y saldo para estabilizar, crecer y redistribuir; pareja de la monetaria \vertema{3.A.36}. \emph{Instrumentos}: gasto (consumo, inversión, transferencias), impuestos (suma fija o distorsionantes), deuda. \emph{Discrecional} vs.\ \emph{automática} (\concepto{estabilizadores automáticos} \vertema{4.B.4}). 3 funciones de \autor{Musgrave} (1959): asignativa, redistributiva, estabilizadora \vertema{4.B.1}. 2 caminos: \emph{corto plazo} (\autor{Kahn}, 1931 $\to$ ER, \autor{Barro} 1974 $\to$ RBC, \autor{Baxter}-\autor{King} 1993 $\to$ NEK, \autor{Woodford} 2011 $\to$ \autor{Blanchard}-\autor{Leigh} 2013) y \emph{largo} (\autor{Diamond} 1965; \autor{Chamley} 1986-\autor{Judd} 1985; \autor{Barro} 1990, \autor{Rebelo} 1991).
\item \textbf{Problemática}: ¿cuánto mueve $Y$ el gasto a CP y de qué depende? ¿Qué hacen deuda, impuestos y pensiones con el ahorro y el capital? ¿Cambia la \emph{tasa} de crecimiento o solo el \emph{nivel}?
\item \textbf{Estructura}: \hyperref[rep:I]{I}, corto plazo (IS-LM, RBC, NEK); \hyperref[rep:II]{II}, ahorro (horizonte infinito; generaciones solapadas); \hyperref[rep:III]{III}, crecimiento (gasto productivo, impuestos). Del CP al LP y del \emph{nivel} a la \emph{tasa}: ahorro antes que crecimiento, aunque el título los nombre al revés.
\end{itemize}

% ==========================================================================
\bloque[9,5]{I. Política fiscal en el corto plazo\label{rep:I}}
\begin{ideaclave}
El gasto mueve la producción a CP, pero cuánto depende del modelo: mucho en el keynesiano (precios rígidos); poco en el RBC (efecto riqueza, multiplicador $>0$ y $<1$); algo más en el NEK, según responda el banco central. En todos el gasto se paga, hoy o mañana: eso lleva al largo plazo.
\end{ideaclave}

\sub{I.1. Debate tradicional: IS-LM\label{rep:I.1}}
\begin{itemize}
\item \emph{Idea}: precios rígidos $\Rightarrow$ la demanda fija $Y$; \concepto{multiplicador} de \autor{Kahn} (1931) por el consumo inducido; \autor{Keynes} prefería gasto a rebajas (parte se ahorra) \vertema{3.A.4}. Con $C=\overline C+c(1-t)Y$, $I=\overline I-bi$, $M/P=kY-hi$ \vertema{3.A.5}:
\begin{formula}
\centering $\displaystyle\frac{\Delta Y^*}{\Delta\overline G}=\frac{1}{1-c(1-t)+bk/h}$
\end{formula}
\item[] $bk/h$ (se \emph{suma}) = \concepto{expulsión} (\emph{crowding-out}): $G$↑ $\Rightarrow$ $Y$↑ $\Rightarrow$ $M^d$↑ $\Rightarrow$ $i$↑ $\Rightarrow$ $I$↓. \emph{Escuelas} ($b$, $h$): \emph{clásicos} $h=0$, LM vertical, expulsión total; \emph{keynesianos} $b$ pequeño, $h$ grande (trampa de liquidez): $\approx1/[1-c(1-t)]$; \emph{monetaristas} $b$ grande, $h$ pequeño: expulsión casi total (\autor{Friedman}-\autor{Heller}, 1969; LP: \autor{Phillips} vertical, solo inflación \vertema{3.A.28}). Presupuesto equilibrado: multiplicador 1 (\autor{Haavelmo}, 1945), $<1$ en IS-LM.
\item \emph{Abierta} (\autor{Mundell}-\autor{Fleming}): $m$ reduce el multiplicador; movilidad perfecta: muy eficaz con tipo \emph{fijo}, ineficaz con \emph{flexible} (apreciación); sin movilidad, al revés \vertema{3.A.32}.
\item \emph{Valoración}: estático, sin expectativas ni restricción intertemporal; desde los 70: ¿qué hace $G$ si los hogares saben que lo pagarán?
\end{itemize}

\sub{I.2. Modelo de ciclo real con política fiscal\label{rep:I.2}}
\begin{itemize}
\item \emph{Idea}: RBC \vertema{3.A.42} + Estado que gasta y grava (\autor{Christiano}-\autor{Eichenbaum}, 1992; \autor{Baxter}-\autor{King}, 1993; presentación de \autor{Sims}, 2017). ¿Qué hace $G$↑ sin rigideces y con hogares que saben que lo pagarán?
\item \emph{Supuestos}: los del RBC (cerrada, 1 bien, horizonte infinito, representativos, competencia perfecta, precios flexibles, ER, \emph{Cobb-Douglas}, $\ln A_t$ AR(1)). $G_t$ no da utilidad ni produce; $\ln G_t=(1-\rho_G)\ln\overline G+\rho_G\ln G_{t-1}+\varepsilon^G_t$. Financia con $T_t$ (suma fija), $\tau_{L,t}$, $\tau_{K,t}$ y deuda:
\begin{formula}
\centering $G_t+r^D_{t-1}D_{t-1}=T_t+\tau_{L,t}w_tL_t+\tau_{K,t}r_tK_t+(D_t-D_{t-1})$
\end{formula}
\item \emph{Ecuaciones}: las del RBC con 2 cambios: cuñas fiscales en el hogar y $G$ en la restricción de recursos:
\begin{formula}
\centering $C_t^{-\theta}=\beta\,\mathbb{E}_t\big[C_{t+1}^{-\theta}\big(1+(1-\tau_{K,t+1})r_{t+1}\big)\big]$\\
$\chi L_t^{\varphi}=C_t^{-\theta}(1-\tau_{L,t})\,w_t$
\end{formula}
\funcionclave{K_{t+1}=(1-\delta)K_t+Y_t-C_t-G_t}
\item[] La restricción del hogar sobra (\autor{Walras}).
\item \concepto{Equivalencia ricardiana}: tabla de 8 condiciones (tecnología, $w$, $r$, \autor{Euler}, oferta de trabajo, bonos, recursos, gasto): \emph{¡ninguna depende de $D_t$ o de $T_t$!} Dada la senda de $G$, da igual impuestos o deuda (\autor{Barro}, 1974; \autor{Ricardo}, 1817, que dudó; nombre de \autor{Buchanan}, 1976). \emph{La deuda no es riqueza neta, sino impuestos aplazados}: no hay plazo que no se cumpla ni deuda que no se pague. \emph{Supuestos que pueden fallar}: horizonte infinito o herencias (\autor{Blanchard}, 1985), suma fija, mercados perfectos (liquidez, tipo del Estado menor), certidumbre, racionalidad \vertema{3.A.33}, senda de $G$ dada. \emph{No keynesianos} (\autor{Giavazzi}-\autor{Pagano}, 1990: Dinamarca 1983-86, Irlanda 1987-89): riqueza (menos impuestos futuros), prima de riesgo, incertidumbre (\autor{Bertola}-\autor{Drazen}, 1993), oferta (salarios públicos; \autor{Alesina}-\autor{Perotti}, 1995); recortar gasto cuesta menos que subir impuestos (\autor{Alesina}-\autor{Ardagna}, 2010), pero la mayoría de las consolidaciones contrae a CP (\autor{Guajardo}-\autor{Leigh}-\autor{Pescatori}, 2014) \vertema{3.A.39}.
\item \emph{Gasto↑ con suma fija}: \concepto{efecto riqueza negativo} $\Rightarrow$ $C$↓ y, ocio normal, $L^s$ a la dcha.: $w$↓, $L$↑, $Y$↑ (en el impacto, $r$↑: sustitución intertemporal). $Y$ sube menos que $G$ (expulsa $C$ e $I$): multiplicador \emph{positivo y menor que 1}.
\item \emph{Persistencia}: \emph{transitorio} ($\rho_G=0$): apenas riqueza, se paga con $I$↓; impacto 0,03. \emph{Persistente} ($\rho_G=0{,}95$): 0,32 en el impacto, $C$ cae más. \emph{Permanente}: $K$ sube como $L$, multiplicador LP 0,75 ($1/[s_C(\theta+\varphi)/\theta+s_G]$); $>1$ con trabajo más elástico. Oferta inelástica $\Rightarrow$ \concepto{expulsión completa}.
\end{itemize}
\graficorepaso[\linewidth]{irf-rbc-gasto}{Gasto +1\,\% del PIB; $\alpha=1/3$, $\beta=0{,}99$, $\delta=0{,}025$, $\theta=\varphi=1$, $G/Y=0{,}2$. a) $Y$, $C$, $I$ con suma fija, $\rho_G=0{,}95$ (continuas) y $0$ (discontinuas); b) horas↑, $w$↓; c) $Y$ con suma fija (0,32) y con impuesto sobre el trabajo ($-0{,}68$)}
\begin{itemize}
\item \emph{Impacto} de +1\,\% del PIB ($\Delta Y$, $\Delta C$, $\Delta I$; horas en \%): suma fija con $\rho_G=0$: 0,03, $-0{,}04$, $-0{,}93$; 0,05. Con $\rho_G=0{,}95$: 0,32, $-0{,}36$, $-0{,}31$; 0,49. Impuesto sobre el trabajo: $-0{,}68$, $-0{,}61$, $-1{,}07$; $-1{,}02$.
\item \emph{Extensiones}: \emph{distorsionantes}: la ER falla; $\tau_L$↓ salario neto y horas: sustitución domina, multiplicador \emph{negativo} (tipo implícito sobre el trabajo 2024: 36,2\,\% España, 37,1\,\% UE); $\tau_K$↓ inversión; tipos como perturbación (\autor{McGrattan}, 1994). \concepto{Suavización impositiva} (\autor{Barro}, 1979): distorsión $\propto\tau^2$ $\Rightarrow$ tipo medio constante $T_t/Y_t=\bar\tau$, la deuda absorbe lo transitorio (guerra, pandemia) y el ciclo: \autor{Ramsey} en el tiempo \vertema{4.B.11}. \emph{Gasto productivo} (\autor{Baxter}-\autor{King}): $Y=AK^{\alpha}L^{1-\alpha}(K^G)^{\alpha_G}$; 3 efectos: absorbe recursos, $Y$↑ con los mismos factores, PMg privados↑ (complementa); CP domina el 1.º, LP el 3.º; con $\alpha_G=1-\alpha$, crecimiento endógeno.
\item \emph{Evidencia}: acierta en $Y$ y horas, no en $C$: en los VAR, $C$ \emph{sube} (\autor{Blanchard}-\autor{Perotti}, 2002; \autor{Galí}-\autor{López-Salido}-\autor{Vallés}, 2007).
\item \emph{Valoración}: \emph{a favor}, microfundamentos (efecto riqueza): importan el gasto y su financiación, no el déficit. \emph{En contra}: $C$ cae; necesita oferta de trabajo más elástica que la micro \vertema{3.A.25}; ciclo eficiente: nada que estabilizar.
\end{itemize}

\sub{I.3. Modelo neokeynesiano\label{rep:I.3}}
\begin{itemize}
\item \emph{Supuestos}: canónico \vertema{3.A.7} sin capital, competencia monopolística, \autor{Calvo}; $Y_t=C_t+G_t$, suma fija; $g_t\equiv-\ln(1-G_t/Y_t)$, $c_t=y_t-g_t$. El gasto entra en la IS:
\funcionclave{y_t=\mathbb{E}_ty_{t+1}-\tfrac1\theta\big(i_t-\mathbb{E}_t\pi_{t+1}-\rho\big)+\big(g_t-\mathbb{E}_tg_{t+1}\big)}
\item[] y en la producción natural, $y^n_t=\psi a_t+\Gamma g_t-\mu/(\theta+\varphi)$, $\Gamma=\theta/(\theta+\varphi)$ (efecto riqueza). \autor{Phillips} NK, $\pi_t=\beta\mathbb{E}_t\pi_{t+1}+\kappa x_t+u_t$, y \autor{Taylor}, $i_t=\rho+\phi_\pi\pi_t+\phi_xx_t+\varepsilon_t$, sin cambios. Solución $x_t=\psi_xg_t$: multiplicador $\Gamma+\psi_x$. Caja: el plano $(x,\pi)$ exige log-linealizar, AR(1), coeficientes indeterminados y expectativas proporcionales.
\item \emph{3 vías}: efecto riqueza (como el RBC); rigidez de precios: \emph{margen contracíclico}, amplifica; respuesta del BC (\autor{Taylor}: $r$↑, expulsa $C$).
\item \emph{El multiplicador depende del BC} (\autor{Woodford}, 2011): $\Delta Y/\Delta G=1+\Delta C/\Delta G$. Flexibles: $\Gamma$ (0,5); \autor{Taylor}: entre $\Gamma$ y 1 (0,87); $r$ constante: \emph{exactamente 1}; \emph{límite inferior}: $\pi^e$↑, $r$↓: $>1$ (1,31 si dura 5 trimestres; 2,25 si 6,7; \autor{Christiano}-\autor{Eichenbaum}-\autor{Rebelo}, 2011; \autor{Eggertsson}, 2011) \vertema{3.A.40}.
\item \emph{Hogares al día}: fracción $\lambda$ (\autor{Campbell}-\autor{Mankiw}, 1989: $\approx$0,5 en EE.\,UU.) $\Rightarrow$ ER falla, $C$↑ con $G$, multiplicador $>1$ (\autor{Galí}-\autor{López-Salido}-\autor{Vallés}, 2007).
\item \emph{Evidencia}: gasto, 0,6-1 (\autor{Ramey}, 2019); \autor{Hall} (2009): 0,7-1 normal, $\approx$1,7 con tipo en cero; FMI supuso $\approx$0,5 y en la crisis $>1$ (\autor{Blanchard}-\autor{Leigh}, 2013). Impuestos: +1\,\% del PIB $\Rightarrow$ $Y$ $\approx-3$\,\% (\autor{Romer}-\autor{Romer}, 2010); $-0{,}78$ a $-1{,}33$ (\autor{Blanchard}-\autor{Perotti}, 2002).
\item \emph{Valoración}: más papel estabilizador y explica multiplicadores altos con tipos en cero; sin capital ni crecimiento: nada del LP.
\end{itemize}

% ==========================================================================
\bloque[9]{II. Efectos a largo plazo sobre el ahorro\label{rep:II}}
\begin{ideaclave}
Con hogares de vida infinita, la deuda no cambia el ahorro nacional, pero un impuesto sobre las rentas del capital sí: a LP conviene no gravarlas. Con generaciones solapadas, el reparto y la deuda reducen ahorro y capital, lo que solo es bueno si la economía acumula de más.
\end{ideaclave}

\sub{II.1. Agente representativo de horizonte infinito\label{rep:II.1}}
\begin{itemize}
\item \emph{Marco}: RCK \vertema{3.A.43} (RBC sin perturbaciones, trabajo inelástico): CRRA, $n$, $g$; $\tilde k$, $\tilde c$ por trabajador eficiente (caja: del RBC al RCK en 4 pasos). Estado: $G$ sin utilidad ni producción; suma fija, deuda o $\tau_K$.
\item \emph{$G$ con suma fija}: \autor{Euler} igual; $\dot{\tilde k}=f(\tilde k)-\tilde c-\tilde G-(n+g+\delta)\tilde k$. \emph{Permanente}: $\dot{\tilde k}=0$ baja en su cuantía, $c$ cae lo mismo al anunciarse; $\tilde k$ y $r$ no cambian: \concepto{expulsión completa del consumo privado}. \emph{Temporal}: $c$ cae menos; mientras dura, $\tilde k$↓ y $r$↑.
\item \emph{Ahorro nacional} $S^N=S^P+S^G$: rebaja con deuda $\Rightarrow$ el hogar ahorra toda la rebaja: $S^P$ sube lo que baja $S^G$, $S^N$, $k$ y $r$ igual (ER) \vertema{3.A.39}. Datos: compensación parcial; España: $S^G<0$ de 2009 a 2024; $S^P$ hogares solo ↑ con fuerza en 2020; 2025: tasa de ahorro de hogares 11,9\,\% vs.\ 14,8\,\% zona del euro.
\item \emph{Impuesto sobre el capital}: entra en \autor{Euler}:
\funcionclave{\frac{\dot{\tilde c}}{\tilde c}=\frac{(1-\tau_K)\big(f'(\tilde k)-\delta\big)-\rho-\theta g}{\theta}\ \Rightarrow\ f'(\tilde k^*)=\delta+\frac{\rho+\theta g}{1-\tau_K}}
\item[] $\dot{\tilde c}=0$ a la izda.: menos $\tilde k^*$, $\tilde y^*$, $\tilde c^*$; al subir, $c$ salta arriba y luego desacumula. Consumo (constante) y trabajo (inelástico): como suma fija.
\item \concepto{Resultado de \autor{Chamley}-\autor{Judd}} (\autor{Judd}, 1985; \autor{Chamley}, 1986): con compromiso y convergencia al estado estacionario, \emph{el impuesto óptimo sobre el capital a LP es nulo}, aunque se quiera redistribuir (\autor{Judd}) y haya que usar distorsionantes (\autor{Chamley}). El capital, gallina de los huevos de oro. \emph{Intuición}: grava el consumo futuro con cuña creciente: $(1+r)^T$ vs.\ $[1+(1-\tau_K)r]^T$ (\autor{Atkeson}-\autor{Chari}-\autor{Kehoe}, 1999): \autor{Ramsey}, no gravar distinto bienes parecidos \vertema{4.B.11}; en el EE la oferta de capital es infinitamente elástica: recae en el trabajo. \emph{Implicación}: gravar el capital instalado y comprometerse a no gravar el futuro: \emph{inconsistencia temporal} \vertema{3.A.36}.
\item \emph{Extensiones}: \autor{Straub}-\autor{Werning} (2020): en \autor{Judd}, con EIS $<1$, óptimo positivo y alto; en \autor{Chamley}, límites activos para siempre. Mercados incompletos (\autor{Aiyagari}, 1995) o vidas finitas: positivo, $\approx$36\,\% (\autor{Conesa}-\autor{Kitao}-\autor{Krueger}, 2009). \emph{Evidencia}: tipo implícito sobre el capital 2024: 29,5\,\% España, 32,5\,\% UE.
\item \emph{Valoración}: con crecimiento exógeno solo cambia el \emph{nivel}, nunca la tasa ($g$); y la deuda es neutral: hacen falta generaciones solapadas.
\end{itemize}

\sub{II.2. Pensiones públicas y ahorro\label{rep:II.2}}
\begin{itemize}
\item \emph{Idea}: la mayor transferencia entre generaciones (España: 13,1\,\% del PIB en 2022; 1.ª pensión, \autor{Bismarck}, ley de 1889). \concepto{Capitalización}: lo cotizado se invierte y se devuelve con rendimientos. \concepto{Reparto}: cotizaciones de hoy pagan pensiones de hoy (España, mayoría de Europa) \vertema{4.B.24}. Promesa entre generaciones: «hijo eres, padre serás; cual hiciste, tal habrás». Marco: \autor{Diamond} (1965) \vertema{3.A.29}.
\item \emph{Supuestos}: viven 2 periodos; de jóvenes, 1 unidad de trabajo inelástica, cobran $w_t$, cotizan $\tau$, consumen $C^J_t$ y ahorran $S_t$; de viejos, ahorro + rendimiento + pensión $Z_{t+1}$; $U=\ln C^J_t+\beta\ln C^V_{t+1}$; población a $n$; $Y=K^{\alpha}(AL)^{1-\alpha}$, $A$ a $g$, $\delta=1$ (periodo $\approx$30 años).
\item \emph{Base, sin pensiones}: $\max\ln C^J_t+\beta\ln C^V_{t+1}$ s.a. $C^J_t+S_t=w_t$, $C^V_{t+1}=(1+r_{t+1})S_t$: $S_t=\frac{\beta}{1+\beta}w_t$; $\tilde k_{t+1}=\frac{\beta(1-\alpha)}{(1+\beta)(1+n)(1+g)}\tilde k_t^{\alpha}$. En rojo, lo que añade cada sistema.
\item \emph{Capitalización}: $C^J_t+S_t=\nuevo{(1-\tau)}w_t$, $C^V_{t+1}=(1+r_{t+1})S_t\nuevo{+(1+r_{t+1})\tau w_t}$; $S_t=\frac{\beta}{1+\beta}w_t\nuevo{-\tau w_t}$: ahorro total y capital \emph{no cambian} (obligatorio sustituye 1 a 1 al voluntario, si $\tau$ no excede lo que ahorraría).
\item \emph{Reparto}: $C^J_t+S_t=\nuevo{(1-\tau)}w_t$, $C^V_{t+1}=(1+r_{t+1})S_t\nuevo{+(1+n)\tau w_{t+1}}$:
\funcionclave{S_t=\frac{\beta}{1+\beta}\nuevo{(1-\tau)}w_t\nuevo{-\frac{1}{1+\beta}\,\frac{(1+n)\tau w_{t+1}}{1+r_{t+1}}}}
\item[] $S$↓ por 2 vías (cotización ↓ renta de hoy, pensión ↑ la de mañana); nadie ahorra por los jóvenes. Ley de movimiento (desplazada abajo):
\begin{formula}
\centering $\displaystyle \tilde k_{t+1}=\frac{\beta(1-\alpha)\nuevo{(1-\tau)}}{(1+n)(1+g)\big[1+\beta\nuevo{+\tau(1-\alpha)/\alpha}\big]}\,\tilde k_t^{\alpha}$
\end{formula}
\end{itemize}
\graficorepaso[0.8\linewidth]{diamond-pensiones}{$\alpha=0{,}3$, $\beta=0{,}9$: el reparto baja la ley de movimiento; con $\tau=0{,}045$, $\tilde k^*_{\mathrm R}=\tilde k_{\mathrm{oro}}$}
\begin{itemize}
\item \emph{Rendimientos}: capitalización $1+r$; reparto $(1+n)(1+g)$. \concepto{Condición de \autor{Aaron}} (1966): reparto rinde más si $n+g>r$ = \concepto{ineficiencia dinámica} ($r^*<n+g$) \vertema{3.A.29}; introducir el reparto lleva $\tilde k$ hacia la regla de oro y mejora a \emph{todas} las generaciones (aquí basta $\tau=4{,}5$\,\%). ¿Sin fallos de mercado? \emph{Doble infinitud} (\autor{Shell}, 1971): valor infinito de las dotaciones, falla el 1TFEB \vertema{3.A.22}; y no se puede prestar a quien no ha nacido (\autor{Malinvaud}, 1953; \autor{Cass}, 1972).
\item \emph{Si es eficiente} ($r>n+g$): mal negocio para las futuras; la 1.ª generación cobra sin cotizar; volver a capitalización: una generación paga 2 veces.
\item \emph{La deuda hace lo mismo}: $K_{t+1}=L_tS_t-D_t$: desplaza capital y traslada la carga a los no nacidos (\autor{Diamond}, 1965); ER solo con herencias altruistas (\autor{Barro}, 1974).
\item \emph{Extensiones}: jubilación endógena: efecto ambiguo (\autor{Feldstein}, 1974); muchas generaciones (\autor{Auerbach}-\autor{Kotlikoff}, 1987); \concepto{paternalismo libertario} (\autor{Thaler}-\autor{Sunstein}, 2008) \vertema{3.A.1}: R.\,U., afiliación automática desde oct.\ 2012; Dinamarca: 1 corona de subvención fiscal crea 1 céntimo de ahorro nuevo; las aportaciones automáticas, mucho (\autor{Chetty} et al., 2014).
\item \emph{Evidencia}: \autor{Feldstein} (1974): la riqueza de la Seguridad Social reduce mucho el ahorro; \autor{Leimer}-\autor{Lesnoy} (1982): error de programación, efecto menor y poco robusto. Países con reparto ahorran menos; los que pasaron a aportación definida no ahorran más después (\autor{Samwick}, 2000). Beneficios del capital > inversión: no hay ineficiencia dinámica (\autor{Abel} et al., 1989): el reparto es transferencia a las 1.as generaciones, no mejora de \autor{Pareto}. \emph{Envejecimiento}: $n$↓ $\Rightarrow$ rendimiento del reparto↓. España: mayores de 65/20-64 años, 33,3\,\% (2022) $\to$ 64,5\,\% (2054); pensiones 13,1\,\% del PIB (2022) $\to$ máx.\ 17,3\,\% ($\approx$2050) $\to$ 16,7\,\% (2070); UE 11,4\,\% $\to$ máx.\ 12,2\,\% (2035-44) $\to$ 11,8\,\%.
\item \emph{Valoración}: deja fuera el riesgo (mercados vs.\ demografía y política), la jubilación endógena y la redistribución intrageneracional. Reformas mixtas: cuentas nocionales (Suecia), fondos de reserva, MEI español \vertema{4.B.24}.
\end{itemize}

% ==========================================================================
\bloque[8]{III. Efectos a largo plazo sobre el crecimiento\label{rep:III}}
\begin{ideaclave}
Con rendimientos decrecientes, la política fiscal solo cambia el nivel; si el capital no los tiene, también la tasa. Con gasto productivo, el tamaño del Estado tiene un óptimo, $G/Y=1-\alpha$, y los impuestos que gravan la acumulación reducen la tasa de crecimiento. Lo que más importa es la composición.
\end{ideaclave}

\sub{III.1. Crecimiento endógeno: \autor{Barro} (1990)\label{rep:III.1}}
\begin{itemize}
\item \emph{Idea}: 2 caras del Estado: suministra bienes productivos (infraestructuras, justicia, seguridad, I+D), pero los financia con impuestos que reducen el rendimiento de invertir: «ni contigo ni sin ti». Si el gasto fuera improductivo, tamaño óptimo nulo.
\item \emph{Supuestos}: $Y=AK^{\alpha}G^{1-\alpha}$, $G$ \emph{flujo} que las empresas toman como dado; $\tau$ proporcional sobre la renta; $G=\tau Y$ (presupuesto equilibrado); ahorro constante $s$ (\autor{Solow}), población a $n$.
\item \emph{Desarrollo}: $\dot k=s(1-\tau)y-(n+\delta)k$; restricción del Estado: $G/L=(\tau A)^{1/\alpha}k$; $y=A^{1/\alpha}\tau^{(1-\alpha)/\alpha}k$: un $AK$, sin transición \vertema{3.A.44}:
\begin{formula}
\centering $\displaystyle\gamma=s\,(1-\tau)\,A^{1/\alpha}\tau^{(1-\alpha)/\alpha}-(n+\delta)$
\end{formula}
\item[] \emph{2 efectos}: \textcolor{tceerojo}{$(1-\tau)$ ($-$, menos renta disponible)} y \textcolor{tceeverde}{$\tau^{(1-\alpha)/\alpha}$ (+, capital más productivo)}. Maximizando:
\funcionclave{\tau^*=\frac{G}{Y}=1-\alpha}
\item[] U invertida; \concepto{regla de \autor{Barro}}: $\partial Y/\partial G=(1-\alpha)Y/G=1$; $\tau^*$ = elasticidad de $Y$ al gasto. Con \autor{Ramsey}: $\gamma=[(1-\tau)\alpha A^{1/\alpha}\tau^{(1-\alpha)/\alpha}-\delta-\rho]/\theta$, mismo $\tau^*$.
\end{itemize}
\graficorepaso[0.85\linewidth]{barro}{$\gamma(\tau)$: efecto del gasto (+) a la izquierda de $\tau^*=1-\alpha$, del impuesto ($-$) a la derecha; en $\tau=0$ y $\tau=1$, $\gamma=-(n+\delta)$}
\begin{itemize}
\item \emph{Implicaciones}: el tamaño del Estado afecta al crecimiento a LP; la \emph{composición} importa (improductivo: solo el efecto negativo; impuesto que no grave invertir no frena). \emph{Equilibrio no óptimo}: el hogar recibe solo $\alpha$ de lo que rinde su capital; el planificador crece a $\gamma^{\mathrm P}=[\alpha A^{1/\alpha}(1-\alpha)^{(1-\alpha)/\alpha}-\delta-\rho]/\theta$; suma fija alcanzaría el óptimo (\autor{Barro}-\autor{Sala-i-Martin}, 2004). Matices: rinde más con reformas de oferta; efectos sobre todo a LP.
\item \emph{Extensiones}: \emph{congestión}: importa $G/Y$, el impuesto actúa como precio (\autor{Barro}-\autor{Sala-i-Martin}, 1992); \emph{capital público acumulable}: transición; RCE de $K$ y $K^G$ juntos.
\item \emph{Valoración}: base al tamaño óptimo del gasto productivo; pero filo de la navaja del $AK$, gasto homogéneo sin deuda, elasticidad difícil de medir.
\end{itemize}

\sub{III.2. Estructura impositiva y crecimiento\label{rep:III.2}}
\begin{itemize}
\item \emph{Teoría}: $AK$ (\autor{Rebelo}, 1991) \vertema{3.A.44}: $\tau_K$ reduce la \emph{tasa} para siempre (en RCK, solo el nivel: \hyperref[rep:II.1]{II.1}):
\begin{formula}
\centering $\displaystyle\gamma=\frac{(1-\tau_K)A-\delta-\rho}{\theta}$
\end{formula}
\item[] $\Rightarrow$ mejor gravar el consumo que el capital o el trabajo (aunque el IVA distorsiona consumo-ocio con trabajo elástico) \vertema{4.B.11}.
\item \emph{Jerarquía de la OCDE} (\autor{Johansson} et al., 2008; \autor{Arnold}, 2011), de más a menos dañino: sociedades $>$ renta personal $>$ consumo $>$ propiedad inmobiliaria recurrente.
\item \emph{Gasto productivo}: \autor{Aschauer} (1989): elasticidad al capital público $\approx$0,4; metaanálisis \autor{Bom}-\autor{Ligthart} (2014, 578 estimaciones, 68 estudios): 0,083 CP y 0,122 LP. España: inversión pública 5,2\,\% (2009) $\to$ 2,0\,\% (2016-17) $\to$ 3,0\,\% (2025) vs.\ 3,9\,\% UE \vertema{3.A.34}.
\item \emph{Restricción completa}: \autor{Kneller}-\autor{Bleaney}-\autor{Gemmell} (1999, OCDE): distorsionantes ↓ crecimiento, gasto productivo ↑, improductivo neutro; ojo con la causalidad inversa \vertema{3.A.45}.
\item \emph{Deuda}: umbral del 90\,\% (\autor{Reinhart}-\autor{Rogoff}, 2010) no resiste la réplica (\autor{Herndon}-\autor{Ash}-\autor{Pollin}, 2014: errores, sin umbral); causalidad en los 2 sentidos \vertema{3.A.39}.
\end{itemize}

% ==========================================================================
\bloque[1]{Conclusión}
\begin{itemize}
\item \textbf{Recapitulación}: «A modo de conclusión,» CP: multiplicador $>0$ y $<1$ con hogares ricardianos y precios flexibles; mayor con rigideces y BC que acomoda. Ahorro: horizonte infinito, deuda neutral, $\tau_K$ reduce $k$ (\autor{Chamley}-\autor{Judd}); generaciones solapadas, reparto y deuda ↓ ahorro y capital, bueno solo con ineficiencia dinámica. Crecimiento: óptimo $G/Y=1-\alpha$; gravar la acumulación baja la \emph{tasa}.
\item \textbf{Relevancia}: deudas $\approx$100\,\% y envejecimiento: en qué gastar y cómo financiarlo. BCE: facilidad de depósito 2\,\% (jun.\ 2025), 2,25\,\% (jun.\ 2026), 2,50\,\% (sept.\ 2026): margen fiscal según deuda heredada y $r$ vs.\ $g$ (\autor{Blanchard}, 2019) \vertema{3.A.39}.
\item \textbf{Extensiones}: \emph{policy mix} y quién cierra la restricción \vertema{3.A.36} \vertema{3.A.39}; multiplicadores en el límite inferior \vertema{3.A.40}; imposición óptima \vertema{4.B.11}; pensiones en España \vertema{4.B.24}; evidencia del crecimiento \vertema{3.A.45}.
\item \textbf{Opinión}: se juzgó por el multiplicador a CP, pero lo duradero es la composición: estabilizadores automáticos en el ciclo (suavización) y, a LP, proteger la inversión productiva y gravar poco el ahorro y la acumulación.
\item \textbf{Idea final}: «quien no trabaja en verano, pasa hambre en invierno»: la política fiscal puede hacer de hormiga. En invierno, gastar lo guardado (sostener la demanda); a LP, gastar en lo que rinde y no castigar a quien ahorra: el ahorro decide el \emph{nivel}; el gasto productivo y los impuestos que no gravan la acumulación, el \emph{ritmo}; y unas pensiones sostenibles, el invierno de cada generación.
\end{itemize}

% ==========================================================================
\bloque{Autores y fechas}
\hito{1817}{\autor{Ricardo}: cómo pagar una guerra (cap.~17)}
\hito{1889}{\autor{Bismarck}: 1.ª pensión pública de vejez}
\hito{1931}{\autor{Kahn}: el multiplicador}
\hito{1945}{\autor{Haavelmo}: presupuesto equilibrado}
\hito{1959}{\autor{Musgrave}: las 3 funciones}
\hito{1965}{\autor{Diamond}: generaciones solapadas con capital}
\hito{1966}{\autor{Aaron}: paradoja del seguro social}
\hito{1971}{\autor{Shell}: doble infinitud}
\hito{1974}{\autor{Barro}: «Are Government Bonds Net Wealth?»; \autor{Feldstein}: Seguridad Social y ahorro}
\hito{1976}{\autor{Buchanan}: nombre de la equivalencia ricardiana}
\hito{1979}{\autor{Barro}: suavización impositiva}
\hito{1982}{\autor{Leimer}-\autor{Lesnoy}: el error de \autor{Feldstein}}
\hito{1985}{\autor{Judd}; 1986, \autor{Chamley}: $\tau_K=0$ a LP}
\hito{1989}{\autor{Aschauer}: capital público; \autor{Campbell}-\autor{Mankiw}; \autor{Abel} et al.}
\hito{1990}{\autor{Barro}: gasto productivo; \autor{Giavazzi}-\autor{Pagano}: no keynesianos}
\hito{1991}{\autor{Rebelo}: $AK$}
\hito{1992}{\autor{Christiano}-\autor{Eichenbaum}; 1993, \autor{Baxter}-\autor{King}: gasto en el RBC}
\hito{1995}{\autor{Aiyagari}: mercados incompletos}
\hito{1999}{\autor{Kneller}-\autor{Bleaney}-\autor{Gemmell}; \autor{Atkeson}-\autor{Chari}-\autor{Kehoe}}
\hito{2002}{\autor{Blanchard}-\autor{Perotti}: VAR fiscal}
\hito{2007}{\autor{Galí}-\autor{López-Salido}-\autor{Vallés}: hogares al día}
\hito{2008}{\autor{Thaler}-\autor{Sunstein}; \autor{Galí}-\autor{Monacelli}: unión monetaria}
\hito{2010}{\autor{Romer}-\autor{Romer}: impuestos; \autor{Reinhart}-\autor{Rogoff}: 90\,\%}
\hito{2011}{\autor{Woodford}; \autor{Christiano}-\autor{Eichenbaum}-\autor{Rebelo}}
\hito{2013}{\autor{Blanchard}-\autor{Leigh}: multiplicadores del FMI}
\hito{2014}{\autor{Bom}-\autor{Ligthart}; \autor{Herndon}-\autor{Ash}-\autor{Pollin}}
\hito{2019}{\autor{Ramey}; \autor{Alesina}-\autor{Favero}-\autor{Giavazzi}}
\hito{2020}{\autor{Straub}-\autor{Werning}: contra $\tau_K=0$}

% ==========================================================================
\bloque{Datos clave}
\begin{itemize}
\item España 2025: déficit \dato{2,4\,\% del PIB}{Eurostat}; deuda \dato{100,5\,\%}{Eurostat}; ahorro nacional bruto \dato{25,0\,\%}{Eurostat, máx.\ desde 1995}; ahorro AA.\,PP.\ \dato{0,3\,\%}{Eurostat}; tasa de ahorro de hogares \dato{11,9\,\%}{Eurostat; 14,8\,\% zona del euro}.
\item Gasto AA.\,PP.\ 2025: \dato{45,3\,\% del PIB}{Eurostat, España}; \dato{49,8\,\%}{Eurostat, zona del euro}.
\item Pensiones: \dato{13,1\,\% (2022) $\to$ 17,3\,\% ($\approx$2050)}{CE, Ageing Report 2024}; dependencia \dato{33,3 $\to$ 64,5\,\%}{CE, 2022-2054}.
\item Tipo implícito 2024: trabajo \dato{36,2\,\%}{CE; UE 37,1}; capital \dato{29,5\,\%}{CE; UE 32,5}.
\item Inversión pública: \dato{5,2\,\% (2009), 2,0\,\% (2016-17), 3,0\,\% (2025)}{Eurostat; UE 3,9\,\%}.
\item BCE, facilidad de depósito: \dato{2,50\,\%}{BCE, sept.\ 2026}.
\end{itemize}

% ==========================================================================
\bloque{Preguntas del tribunal}
\begin{enumerate}
\item \textbf{¿Por qué el multiplicador es menor que 1 en el RBC?} Actúa solo por el efecto riqueza: $L$↑, pero $G$ expulsa $C$ e $I$; mayor cuanto más persistente el gasto y más elástica la oferta de trabajo.
\item \textbf{¿Cuándo supera 1 en el NEK?} Si el BC no sube el tipo real (con $r$ constante, 1), en el límite inferior ($\pi^e$↑, $r$↓) o con hogares que viven al día.
\item \textbf{¿Por qué no gravar el capital a LP si el Estado quiere redistribuir?} Cuña creciente sobre el consumo futuro y oferta de capital infinitamente elástica en el EE: recae en el trabajo (\autor{Chamley}-\autor{Judd}); con mercados incompletos o EIS $<1$, positivo.
\item \textbf{¿Puede el reparto mejorar a todas las generaciones?} Sí, si hay ineficiencia dinámica ($n+g>r$, \autor{Aaron}); los datos (\autor{Abel} et al., 1989) no la encuentran: es transferencia a las 1.as generaciones.
\item \textbf{¿Por qué $\tau_K$ baja el nivel en RCK y la tasa en \autor{Rebelo}?} Con rendimientos decrecientes el PMg del capital se recupera al bajar $k$; en el $AK$ es constante y $\tau_K$ reduce el rendimiento neto para siempre.
\item \textbf{¿Y en economía abierta con HER?} (\autor{Corsetti}-\autor{Pesenti}, 2001; HER = hipótesis de expectativas racionales): con suma fija, ER; a CP (salarios fijados) $\Delta Y=\Delta G$, $\Delta C=\Delta C^*=0$, sin apreciación; a LP mejora la relación real de intercambio y empobrece al vecino (\emph{beggar-thy-neighbour}). En \autor{Mundell}-\autor{Fleming}: muy eficaz con tipo fijo y movilidad perfecta (el BC compra divisas, LM a la dcha.); ineficaz con flexible (apreciación).
\item \textbf{¿Qué falla en el umbral del 90\,\%?} Errores de cálculo y ponderación (\autor{Herndon}-\autor{Ash}-\autor{Pollin}, 2014); causalidad en los 2 sentidos.
\end{enumerate}

\begin{trampa}
\begin{itemize}
\item ER: 11 de 14 preguntas oficiales. Suma fija, mercados perfectos, horizonte infinito (o finito con herencias altruistas): rebaja con deuda $\Rightarrow$ $S^P$↑ lo mismo, $C$ igual; Estado que se endeuda más barato que los hogares la rompe.
\item Junio 2023 (P47): expansión contractiva si se ve permanente y $C$ cae más que sube $G$; con suma fija, $C$ cae como mucho lo mismo: hace falta, p.\,ej., impuesto distorsionante.
\item Multiplicador de \autor{Kahn} (no \autor{Keynes} ni \autor{Hahn}); $bk/h$ se \emph{suma} en el denominador.
\item \autor{Musgrave}: 1959, no 1939. Divina coincidencia: \autor{Blanchard}-\autor{Galí} (2007; DT de 2005).
\item Unión monetaria: \autor{Galí}-\autor{Monacelli} 2008 (la de 2005 es monetaria en abierta).
\item Precios rígidos: margen \emph{contracíclico}.
\item \autor{Aaron}: $n+g>r$, no $r<n$.
\item \autor{Laffer} (nov.\ 2016) va con imposición \vertema{4.B.11}.
\item \autor{Barro} (1990): $G$ es flujo; $\tau^*=1-\alpha$ (elasticidad de $Y$ a $G$), no $\alpha$.
\end{itemize}
\end{trampa}

\end{columnas}
\end{document}
